\section*{Bab \ref{ch:b3:covariant-electromagnetism} --- Elektromagnetisme Kovarian}

\begin{solution}{exo:b3:covariant-electromagnetism:1}
(a) $\gamma = 1.25$: $a'^0 = 1.25(5 - 1.8) = 4.0$, $a'^1 = 1.25(3 -
3) = 0$; $a\cdot a = 25 - 9 = 16 = 16 - 0$. (b) Transformasinya
linear, sehingga ia terdistribusi pada penjumlahan. (c) $(c, \vect v)$: tidak ---
komponen ``waktu''nya sama bagi semua partikel padahal percampurannya
menuntut sebaliknya; $\gamma(c, \vect v) = \dd x^\mu/\dd\tau$: ya, yakni
\omterm{def:b3:covariant-electromagnetism:fourvector}{vektor-empat} yang dibagi invarian $\dd\tau$. (d) Ketiga komponen $\vect E$
bercampur dengan komponen $\vect B$, bukan dengan skalar mana pun: kedua
medannya mengisi tensor berindeks dua, bukan sebuah \omterm{def:b3:covariant-electromagnetism:fourvector}{vektor-empat}.
\end{solution}

\begin{solution}{exo:b3:covariant-electromagnetism:2}
(a) $u = I/nSe = 10/(\num{8.5e28} \times \num{e-6} \times
\num{1.6e-19}) = \qty{7.4e-4}{m/s}$. (b) Kisi: $(nec, \vect 0)$;
elektron: $(-nec, -ne\vect u)$ --- kerapatan arusnya
$\rho_-\vect u = -ne\vect u$ menunjuk berlawanan dengan hanyutannya, dan
itulah \emph{justru} arah arus konvensionalnya. (c) Keduanya statis dan seragam: tiap sukunya lenyap. (d)
Totalnya: $(0, \vect\jmath)$ dengan $j = I/S = \qty{e7}{A/m^2}$ --- yakni
arus murni tanpa muatan, yaitu susunan yang membuat kerelatifan kawat itu
halus.
\end{solution}

\begin{solution}{exo:b3:covariant-electromagnetism:3}
(a) $E_y$: $F^{20} = E_y/c$; $B_x$: $F^{32} = B_x$. (b) Untuk $\vect B
= B\vect e_z$, hanya $F^{12} = -B$ dan $F^{21} = +B$; untuk gelombangnya,
$F^{20} = E/c$, $F^{12} = -E/c$ (beserta pasangan antisetangkupnya). (c)
$F^{12} = \partial^1A^2 - \partial^2A^1 = -\partial_xA_y +
\partial_yA_x = -(\operatorname{\vect{curl}}\vect A)_z = -B_z$: cocok
dengan matriksnya. (d) Laju perubahan energi partikelnya
$\propto F^{0\nu}u_\nu$; dalam kerangka diamnya $u_\nu = (c, \vect 0)$ dan
keantisetangkupannya membuat $F^{00} = 0$: sebuah gaya yang tak pernah
dapat melakukan usaha pada partikel yang diam --- yakni sifat yang
mendefinisikan bagian magnetiknya.
\end{solution}

\begin{solution}{exo:b3:covariant-electromagnetism:4}
(a) Dorongan searah medannya: komponen longitudinalnya tak berubah, $\vect
E' = E_0\vect e_y$, $\vect B' = \vect 0$. (b) Dorongan sepanjang $\vect
e_x$: $E'_y = \gamma E_0$, $B'_z = -\gamma vE_0/c^2$ --- pelatnya
kini mengalir sebagai arus permukaan $\pm\sigma'v$, dan dua lembar arus
berlawanan melingkupi persis medan semacam itu. (c) $E'^2 - c^2B'^2 =
\gamma^2E_0^2(1 - \beta^2) = E_0^2$; $\vect E'\cdot\vect B' = 0$:
keduanya terpelihara. (d) Pelatnya terkerut sepanjang $x$, sehingga $\sigma' =
\gamma\sigma$, selaras dengan $E' = \sigma'/\varepsilon_0 = \gamma
E_0$; jarak antarpelatnya, yang melintang terhadap dorongannya, tak tersentuh.
\end{solution}

\begin{solution}{exo:b3:covariant-electromagnetism:5}
(a) Gaya magnetik $-e\vect v\wedge\vect B$ mendorong elektronnya
sepanjang batangnya: ggl $= vBL$. (b) Dalam kerangka batangnya, batangnya
diam --- tak ada gaya magnetik pada muatan yang diam --- tetapi
transformasinya menyerahkan $\vect E' = \gamma\,\vect v\wedge\vect B$:
sebuah medan listrik sejati yang menjadi penggeraknya. (c) $E'L = \gamma vBL
\approx vBL$ pada orde $v/c$. (d) Ia memang harus bekerja karena ggl
``gerak'' dan ggl ``transformator'' merupakan satu gejala yang dibaca
dalam dua kerangka: aturan fluksnya adalah kovariansi yang mengenakan
pakaian 1831 --- makalah Einstein 1905 dibuka tepat dengan ketaksetangkupan
magnet-dan-penghantar ini.
\end{solution}

\begin{solution}{exo:b3:covariant-electromagnetism:6}
(a) Komponen $x$-nya tak tersentuh; untuk yang melintang,
$E_y'^2 + E_z'^2 - c^2(B_y'^2 + B_z'^2) = \gamma^2[(E_y - vB_z)^2 +
(E_z + vB_y)^2 - c^2(B_y + vE_z/c^2)^2 - c^2(B_z - vE_y/c^2)^2]$;
suku silangnya saling meniadakan dan kuadratnya menghimpun $\gamma^2(1 -
\beta^2) = 1$ kali gabungan yang belum ditransformasikan. (b) Pembukuan
yang sama pada $E_xB_x + E_yB_y + E_zB_z$. (c) $E^2 - c^2B^2 =
3c^2B^2 > 0$: doronglah sepanjang $\vect E\wedge\vect B$ pada $v = c^2B/E =
c/2$; di sana $B' = 0$ dan medan yang bertahan adalah $E' = E/\gamma =
\sqrt3\,cB$ --- yang kuadratnya sama dengan invariannya, sebagaimana seharusnya. (d) Gelombang cahaya mempunyai kedua invariannya
nol: kerangka mana pun harus menghasilkan kembali $E' = cB' \perp$, tak
pernah medan murni.
\end{solution}

\begin{solution}{exo:b3:covariant-electromagnetism:7}
(a) Dengan $v_{\text{d}} = E/B$: $E'_y = \gamma(E - v_{\text{d}}B) =
0$; $B'_z = \gamma(B - v_{\text{d}}E/c^2) = B/\gamma$: medan magnetik
murni yang sedikit melemah. (b) Di sana: sebuah lingkaran pada
$qB'/\gamma m$; kembali di laboratorium: lingkaran itu ditambah hanyutan
seragamnya --- yakni sikloid yang merayap pada $E/B$ tegak lurus terhadap
kedua medannya, yaitu lintasan muatan dalam magnetron. (c) Partikel yang
bergerak tepat pada $v_{\text{d}}$ diam dalam kerangka hanyutnya, tempat
satu-satunya medan bersifat magnetik dan ia tak merasakan apa pun:
tanpa pembelokan. (d) Untuk $E > cB$ hanyutan yang dituntutnya melebihi
$c$: tak ada kerangka semacam itu; sebaliknya kerangka ber-$E$ murni ada
(latihan sebelumnya) dan muatannya dipercepat tanpa batas sepanjang
medannya.
\end{solution}

\begin{solution}{exo:b3:covariant-electromagnetism:8}
(a) Fasenya $\omega t - \vect k\cdot\vect r$ menghitung peristiwa
perlintasan puncak gelombang, yakni bilangan invarian; ia sama dengan
$k^\mu x_\mu$, dan keinvarianan hasil kalinya untuk semua $x^\mu$
memaksa $k^\mu$ bertransformasi secara \omterm{def:b3:covariant-electromagnetism:fourvector}{vektor-empat}. (b) $\omega'/c = \gamma(\omega/c -
\beta k_x)$ dengan $k_x = (\omega/c)\cos\theta$; untuk $\theta = 0$,
$\omega' = \omega\sqrt{(1-\beta)/(1+\beta)}$. (c) $\cos\theta' =
k'_x/(\omega'/c)$: yakni rumus yang dinyatakan itu. (d) $\beta = \num{e-4}$:
$20.5''$; sepanjang setahun kedudukan semunya menyapu elips dengan
setengah sumbu sudut sebesar itu --- yakni aberasi Bradley.
\end{solution}

\begin{solution}{exo:b3:covariant-electromagnetism:9}
(a) Dalam kerangka diamnya, Coulomb; dorongannya mengalikan komponen
melintangnya dengan $\gamma$ (pada bidang melintangnya, $E'_\perp = \gamma
q/4\pi\varepsilon_0b^2$) sedangkan medan searah geraknya, yang dihitung
di depan atau di belakangnya pada jarak laboratorium $r$, memetakan ke
jarak kerangka diam $\gamma r$: berkurang sebesar $1/\gamma^2$. (b)
Cakramnya yang selebar sudut $1/\gamma$ menyapu lewat pada $v$: durasinya
$\sim (b/\gamma)/v$. (c)
$E' = 7250 \times \num{1.6e-19} \times \num{9e9}/\num{e-4} \approx
\qty{0.1}{V/m}$, yang berlangsung $2b/\gamma c \approx \qty{9}{fs}$. (d)
$B' = vE'/c^2 \approx E'/c$ dan kedua invariannya berorde $O(1/\gamma^2)$:
secara setempat tak terbedakan dari pulsa cahaya --- itulah dasar
pemerian ``foton setara'' bagi tumbukan muatan cepat.
\end{solution}

\begin{solution}{exo:b3:covariant-electromagnetism:10}
(a) $E'_y = \gamma(E - vB) = \gamma E_0(1 - \beta)\cos(\cdots) =
E_0\sqrt{(1-\beta)/(1+\beta)}\cos(\cdots)$, dan $B'_z = E'_y/c$ lewat
aljabar yang sama. (b) Fasenya bertransformasi dengan faktor Doppler
yang sama: gelombang null yang sama, lebih merah dan lebih lemah. (c)
Kerapatan energinya turun sebesar kuadrat faktor Dopplernya; cacah
fotonnya tak berubah --- $h\nu$ tiap fotonlah yang membawa seluruh
faktornya. (d) ``Menumpang di sampingnya'' berarti faktornya $\to 0$
ketika $\beta \to 1$: gelombangnya tak pernah menjadi statis karena
invariannya nol --- tak ada kerangka yang membuat cahaya berhenti, dan di
situlah kerelatifan buku ini bermula.
\end{solution}

\begin{solution}{exo:b3:covariant-electromagnetism:11}
(a) Tiap putaran, $\Delta E = P \times 2\pi r/c$: yakni hasil yang
dinyatakan itu. (b)
$\gamma = \num{e11}/\num{5.11e5} = \num{1.96e5}$: $\Delta E \approx
\qty{2.9}{GeV}$ tiap putaran --- tiga persen energi berkasnya
disuntikkan ulang tiap putaran; klystron LEP adalah pemancar radio
terbesar di Bumi. (c) Proton: $\gamma = 7250$, dengan $\gamma^4$
lebih kecil sebesar $(m_{\text{p}}/m_{\text{e}})^4 \approx \num{e13}$: sekitar
\qty{5}{keV} tiap putaran --- dapat diabaikan. (d) $\gamma^4 = (E/mc^2)^4$:
pada energi yang sama elektron memancar $\num{e13}$ kali lebih banyak;
karena itu penumbuk linear (memancar sekali, bukan tiap putaran) atau
cincin raksasa bagi elektron, sedangkan cincin proton berskala dengan
nyaman sampai \qty{27}{km}.
\end{solution}

\begin{solution}{exo:b3:covariant-electromagnetism:12}
(a) $|\vect p|$ tetap (tanpa usaha); $\dot{\vect p} = q\vect
v\wedge\vect B$ memutar $\vect p$ dengan laju $qvB/p = qB/\gamma m$. (b)
Selisihnya tiap putaran $2\pi(\gamma - 1)$; dengan $\gamma - 1$ tumbuh
secara linear sampai $\Gamma$, selisih totalnya $\approx \pi N\Gamma$;
seperempat periodenya $\pi/2$: $\Gamma \approx 1/2N$. Untuk $N = 100$: $\Gamma =
\num{5e-3}$, $E_k \approx \qty{4.7}{MeV}$ --- skala yang tepat; mesin
sungguhan meregangkannya sampai $\sim\qty{20}{MeV}$ dengan tegangan dee
yang tinggi (putaran lebih sedikit). (c) Sinkrosiklotron: sapukanlah RF-nya
turun sepanjang tiap pulsanya agar mengikuti $qB/\gamma m$. Siklotron
isokron: biarkanlah $B(r)$ tumbuh $\propto\gamma(r)$ sehingga nisbahnya tak
pernah berubah dan berkasnya tetap sinambung. (d) $\gamma = 1 + 590/938 = 1.63$:
medan di tepinya harus melebihi medan pusatnya sebesar faktor itu.
\end{solution}

\begin{solution}{pb:b3:covariant-electromagnetism:1}
\textbf{1.} $u = I/nSe = \qty{7.4e-4}{m/s}$; $u^2/c^2 =
\num{6e-24}$.
\textbf{2.} $\lambda = I/u = nSe = \qty{1.36e4}{C/m}$ --- empat belas
kilocoulomb tiap meter pada tiap arusnya.
\textbf{3.} $B = \mu_0I/2\pi d = \qty{2.0e-4}{T}$; $F = qvB =
\num{e-6} \times \num{e5} \times \num{2e-4} = \qty{2e-5}{N}$, yang
terarah \emph{menuju} kawatnya (arus sejajar saling menarik).
\textbf{4.} Kedua muatan garisnya saling meniadakan secara tepat:
$\lambda_{\text{tot}}
= 0$, tanpa medan pada orde nol.
\textbf{5.} Satu elektron berlebih tiap meter: $E = 2k\lambda_1/d =
\qty{2.9e-7}{V/m}$, dengan gaya $\qty{2.9e-13}{N}$ --- $10^8$ kali
lebih kecil daripada gaya magnetiknya. Kawatnya dapat menyembunyikan
seratus juta elektron liar tiap meter sebelum elektrostatika menyaingi
kemagnetan: menisbahkan gayanya kepada $\vect B$ aman dilakukan.
\textbf{6.} $\qty{2e-5}{N}$ setara dua puluh berat butir pasir: mudah
diukur.
\textbf{7.} Kisi: $-v$. Elektron (yang menghanyut pada $u$ berlawanan
dengan $I$, yakni berlawanan dengan dorongannya): berlaju
$(u + v)/(1 + uv/c^2)$ --- lebih cepat daripada kisinya.
\textbf{8.} $\lambda' = \gamma_v(\lambda_{\text{tot}} - vI/c^2) =
-\gamma_v vI/c^2$.
\textbf{9.} $\lambda' = -\num{e5} \times 10/\num{9e16} =
\qty{-1.1e-11}{C/m}$: sekitar $7 \times 10^7$ elektron berlebih tiap
meter. Arus elektronnya, yang lebih cepat dalam kerangka ini, adalah
yang lebih rapat --- tiap arus terkerut oleh $\gamma$-nya sendiri.
\textbf{10.} $E' = 2k|\lambda'|/d = 2 \times \num{9e9} \times
\num{1.1e-11}/0.01 = \qty{20}{V/m}$, yang menunjuk ke kawatnya; gayanya
$qE' = \qty{2e-5}{N}$, bersifat menarik.
\textbf{11.} Serupa dengan $qvB$ sampai pada faktor $\gamma_v = 1 +
\num{5.6e-8}$: tepat hukum transformasi gaya melintang (gaya-empat),
$F_{\text{diam}} = \gamma F_{\text{lab}}$.
\textbf{12.} Kawatnya tetap mengalirkan arus (yang lebih besar),
sehingga $\vect
B' \neq 0$ --- tetapi muatan kita diam di sini, dan medan magnet hanya
mencengkeram muatan yang bergerak.
\textbf{13.} $|\lambda'|/\lambda = \gamma_v vu/c^2 = \num{8.2e-16}$.
\textbf{14.} Ia mengalikan $\lambda/e \approx \num{8.5e22}$ muatan dasar
tiap meter: $\num{8.2e-16} \times \num{8.5e22}
\approx \num{7e7}$ elektron tiap meter --- makroskopis. Materi begitu
sangat bermuatan sehingga kenetralannya setipis mata pisau; kerelatifan
memiringkan pisau itu.
\textbf{15.} $F/L = \mu_0I_1I_2/2\pi d = \num{2e-7} \times 100/0.01 =
\qty{2e-3}{N/m}$ --- rumus yang dahulu \emph{mendefinisikan} ampere.
\textbf{16.} $B = \mu_0nI = 4\pi\times10^{-7} \times \num{e4} \times
10 = \qty{0.13}{T}$; tekanan magnetiknya $B^2/2\mu_0 \approx
\qty{6300}{Pa} \approx \qty{0.6}{N/cm^2}$: mobil utuh tergantung pada
kerelatifan yang terintegralkan.
\textbf{17.} $c \to \infty$ membunuh $\lambda'$, dan bersamanya gayanya,
bagian magnetik medannya, motor, dinamo dan derek tempat rongsokan:
kemagnetan tidak mempunyai keberadaan yang tak relativistik.
\textbf{18.} Dalam keadaan diam muatannya melihat kawat netral
laboratorium: kedua arusnya saling meniadakan sampai ketelitian
kenetralan materi --- yang dikenal secara percobaan sampai ketelitian
yang menakjubkan --- sehingga tak ada gaya yang bekerja.
\textbf{19.} Bahwa gambaran ``sumber yang diberikan'' itu adalah irisan
satu tensor tunggal menurut satu kerangka: $\vect B$ adalah bagian medan
elektromagnetik yang dinisbahkan gerak seorang pengamat kepada arus,
bukan kepada muatan.
\textbf{20.} Spin elektron: tiap elektron adalah magnet dasar (sebuah
``arus'' kuantum yang hakiki), dan besi adalah materi yang di dalamnya
semuanya searah --- \cref{ch:b3:spin-two-level} dan
\cref{ch:b3:electromagnetism-in-matter} membahasnya.
\textbf{21.} Dalam kerangka diam berkasnya tolakannya murni Coulomb dan
\emph{maksimum}; di laboratorium, tarikan magnetik arus sejajarnya
hampir meniadakannya. Tak ada pertentangan: gaya melintang laboratorium
adalah gaya kerangka diamnya dibagi $\gamma$, dan dinamika yang lebih
lambat (\omterm{prop:b3:relativistic-kinematics:dilation}{pemuluran waktu}) melengkapi pembukuannya.
\textbf{22.} Gaya neto laboratoriumnya $= qE_{\text{lab}} - qvB_{\text{lab}} =
(1 - \beta^2)qE_{\text{lab}} = F_{\text{Coulomb}}/\gamma^2$ menurut
transformasinya; dalam kerangka diamnya faktor itu gamblang:
elektrostatika murni, dengan pemuaiannya disaksikan lewat waktu yang memulur.
\textbf{23.} Vektor Poynting: energinya mengalir melalui \emph{medan} di
sekeliling penghantarnya dengan laju hampir $c$, sedangkan elektronnya
sekadar mengaturnya; tembaga kabelnya adalah pemandu, bukan pipa.
\textbf{24.} $vu/c^2 \approx \num{8e-21}$ --- namun Ørsted melihat
kompasnya berayun di samping sebuah kawat pada 1820: pengali $10^{22}$
muatan sudah bekerja seabad sebelum siapa pun dapat menamainya.
\textbf{25.} Kawat netral \qty{10}{A}, yang dipandang dari
\qty{e5}{m/s}, membawa $\qty{-1.1e-11}{C/m}$: selisih pembukuan
$10^{-16}$ dari \omterm{prop:b3:relativistic-kinematics:contraction}{pengerutan panjang}, dikalikan $10^{23}$ muatan tiap
meter, \emph{itulah} gaya magnetiknya --- seluruh kemagnetan adalah
kerelatifan yang diaudit pada kecepatan siput.
\end{solution}
